%==============================================================================
%  IEEE double-column conference paper -- READY FOR OVERLEAF
%  Compiler: pdfLaTeX.  No .bib file needed (bibliography is inline).
%
%==============================================================================
\documentclass[conference]{IEEEtran}
\usepackage{amsmath,amssymb}
\usepackage{booktabs}
\usepackage{graphicx}
\usepackage[T1]{fontenc}

\begin{document}

\title{Validating White-Box Hallucination Detection Against\\
Surface and Topic Confounds on Multilingual Mu-SHROOM}

\author{
\IEEEauthorblockN{Sourin Ghosh}
\IEEEauthorblockA{
Indian Institute of Technology, Johpur\\ 
M25DE1077\\
m25de1077@iitj.ac.in} 
\and
\IEEEauthorblockN{Dr. Subhash Bhagat}
\IEEEauthorblockA{
Indian Institute of Technology, Jodhpur\\ 
Supervisor\\
\ }
}

\maketitle

\begin{abstract}
White-box hallucination detectors read a language model's internal
activations and are increasingly reported to localize hallucinated spans.
Such probes are, however, easy to fool: a token classifier can score well by
exploiting surface form or topic rather than any hallucination
representation. We test whether a white-box signal is genuine in the
multilingual, span-level setting of Mu-SHROOM (SemEval-2025 Task~3) across
Spanish, Czech, Mandarin, and English, using sparse-autoencoder (SAE)
features over a Gemma-2-2b residual stream. A token probe reaches a
detection AUROC of $0.745$, but a twelve-feature surface-only baseline
recovers about $80\%$ of the above-chance signal---surface form is a large,
under-reported confound. The decisive control is a within-item AUROC, which
ranks tokens inside a single answer where topic, language, and generator are
held constant: here the SAE probe scores $0.818$ against a surface within-item
score of $0.730$, in every language, with non-overlapping error bars. The
signal therefore survives both a surface control and a topic control and is a
genuine token-level hallucination signal. Char-level IoU stays modest
($0.40$ vs.\ a $0.55$ leaderboard reference), framing span extraction as
future work. We also report three negative results---claim-type routing, an
internal-vs-external gap, and proxy retrieval add nothing. The control
methodology, not just the detector, is the contribution.
\end{abstract}

\begin{IEEEkeywords}
hallucination detection, interpretability, sparse autoencoders, confound
analysis, multilingual NLP, Mu-SHROOM.
\end{IEEEkeywords}

%==============================================================================
\section{Introduction}
%==============================================================================
Large language models produce fluent text that is frequently unsupported by
fact. \emph{White-box} detectors attack this by reading the model's own hidden
states, on the premise that a model's internal representation encodes whether a
token is grounded. On Mu-SHROOM, the SemEval-2025 multilingual hallucination
span-detection task, internal-state probes have performed competitively.

The premise has a well-known weakness. A token-level classifier trained on
held-out items can achieve a high area under the ROC curve (AUROC) without
learning anything about hallucination, simply by keying on \emph{surface form}
(token length, punctuation, position) or \emph{topic} (content features that
correlate with which items happen to contain hallucinations in a given corpus).
Without explicit controls, a reported AUROC cannot distinguish a hallucination
representation from these confounds. Recent work has voiced exactly this
skepticism about internal-state methods.

This paper asks a narrow but load-bearing question: in the multilingual,
span-level setting, is the white-box signal \emph{real}? We answer it with two
controls that the literature largely omits:
\begin{itemize}
\item a \textbf{surface baseline}---a twelve-feature, model-free probe that
isolates how much apparent performance is surface form; and
\item a \textbf{within-item AUROC}---ranking tokens \emph{inside} a single
answer, where topic, language, and generator are constant, so the score cannot
be topic or surface confound.
\end{itemize}

Our contributions are: (i)~evidence that surface form is a large,
under-reported confound (it recovers $\approx 80\%$ of the pooled above-chance
signal); (ii)~a decisive within-item control showing a genuine token-level
hallucination signal survives both confounds, multilingually
(within-item AUROC $0.818$); (iii)~a demonstration that interpretable SAE
features match raw residuals at no accuracy cost; and (iv)~three honest
negative results that map the design space.

%==============================================================================
\section{Related Work}
%==============================================================================

\subsection{The Mu-SHROOM shared task}
Mu-SHROOM, SemEval-2025 Task~3~\cite{mushroom}, recast hallucination detection
as multilingual \emph{character-span} identification: given a model output, a
system must mark the substrings that are unsupported, scored by character-level
intersection-over-union (IoU) against soft and hard human labels across more
than ten languages. This span-level, multilingual framing is harder than the
sentence-level binary detection of the original SHROOM task and exposes
methods to typologically diverse scripts. The winning UCSC
system~\cite{ucsc} reached an average IoU near $0.55$ by combining model
signals with external evidence, and a Deloitte entry~\cite{deloitte} reported
residual-stream probes on the same data. Reported leaderboard IoU, however,
is computed against a high FLAG-ALL floor on densely hallucinated answers, so a
strong IoU does not by itself establish that a method has learned a
hallucination-specific signal---the gap this paper targets.

\subsection{Black-box and uncertainty-based detection}
A large body of work detects hallucination without internal access: sampling
consistency (e.g.\ generating multiple answers and measuring disagreement),
token- and sequence-level uncertainty, and self-evaluation prompts. These
methods are model-agnostic but conflate uncertainty with falsehood and degrade
when a model is confidently wrong. We do not pursue them here, but our
confidence ``gap'' ablation is in this spirit and, consistent with their known
brittleness, adds nothing on top of the internal-state probe.

\subsection{Internal-state probing}
The premise that a model's activations encode whether its output is truthful
has accumulated support: linear probes on hidden states separate true from
false statements~\cite{azaria}, Orgad et al.~\cite{orgad} show that such
representations encode more about correctness than the decoded text reveals,
and Ferrando et al.~\cite{ferrando} isolate sparse-autoencoder directions
corresponding to whether a model ``knows'' an entity. Against this optimism, a
skeptical line of work~\cite{skeptic} argues that internal-state probes often
track dataset artifacts---topic, length, formatting---rather than truthfulness,
and that apparent accuracy collapses under distribution shift or proper
controls. This tension is exactly what our surface baseline and within-item
control are designed to resolve: we quantify how much of the signal is
confound, and isolate what remains.

\subsection{Sparse autoencoders and interpretability}
Sparse autoencoders (SAEs) decompose dense activations into an overcomplete
basis of sparse, more nearly monosemantic features, improving on raw
neurons~\cite{monosemantic,cunningham}. Pretrained SAE suites such as Gemma
Scope~\cite{gemmascope} make this practical off-the-shelf for Gemma-2
models~\cite{gemma2}. We use SAE features both as a probe substrate and as an
interpretability lens; notably, we find SAE features match raw residuals on
detection, so interpretability here costs no accuracy.

\subsection{Retrieval-augmented verification}
A complementary external route verifies spans against retrieved evidence.
REFIND~\cite{refind} grounds candidate spans in retrieved documents and flags
unsupported content. We include a REFIND-style retrieval baseline; in our
proxy-model setting (the scorer did not generate the answers) it sits at
chance, which we report as a scoping result rather than a verdict on retrieval
in general.


%==============================================================================
\section{Data and Setup}
%==============================================================================
\textbf{Dataset.} We use the Mu-SHROOM test split for four languages. Item and
token counts are in Table~\ref{tab:data}; overall $556$ items, $66{,}946$
answer tokens, of which $\approx 38\%$ are gold-hallucinated.

\begin{table}[t]
\caption{Mu-SHROOM test split used.}
\label{tab:data}
\centering
\begin{tabular}{lcc}
\toprule
Language & Items & Items with hard labels \\
\midrule
Spanish (es) & 152 & 139 \\
Czech (cs)   & 100 & 87 \\
Mandarin (zh)& 150 & 147 \\
English (en) & 154 & 149 \\
\midrule
Total        & 556 & 462 \\
\bottomrule
\end{tabular}
\end{table}

\textbf{Model and features.} We run \texttt{google/gemma-2-2b} and take the
residual stream after layer~12 (\texttt{resid\_post}). Activations are encoded
with the Gemma Scope SAE \texttt{gemma-scope-2b-pt-res-canonical}
(\texttt{layer\_12/width\_16k/canonical}, $16{,}384$ features). Because the
work was developed on Apple-silicon (MPS), we verified numerical fidelity of
the layer-12 residual against CPU (max absolute difference $0.0000$).

\textbf{Probe.} For each token we form its SAE feature vector and train an
$\ell_2$-regularized logistic regression with balanced class weights. We keep
only features firing on at least $0.5\%$ of training tokens
($\approx 4{,}200$ features) to suppress noise from near-dead features. Splits
are \emph{item-level} ($70/30$) so no item appears in both train and test;
results are averaged over five seeds. The pipeline is sketched in
Fig.~\ref{fig:pipe}.

\begin{figure}[t]
\footnotesize
\begin{verbatim}
   answer tokens  (es / cs / zh / en)
              |
   Gemma-2-2b  resid_post  L12
              |
   Gemma Scope SAE  (16,384 feats)
              |
   freq filter  (>= 0.5% of tokens)
              |
   per-token logistic probe
              |
   token hallucination scores
          /              \
   surface             within-item
   baseline             AUROC
   (form control)      (topic control)
\end{verbatim}
\caption{Detection pipeline and the two confound controls.}
\label{fig:pipe}
\end{figure}

\textbf{Metrics.} We report token-level AUROC in two forms: \emph{pooled}
(rank all test tokens together) and \emph{within-item} (rank tokens inside each
item, then macro-average). For span output we report character-level IoU with a
per-language decision threshold. Floors are FLAG-ALL (mark every character) and
FLAG-NONE (mark nothing).

%==============================================================================
\section{Method: the two controls}
%==============================================================================
\textbf{Surface baseline.} A model-free probe over twelve hand-built features
per token: character length, leading-space flag, punctuation and digit flags,
uppercase ratio, relative position in the answer, and a language one-hot.
Trained identically to the SAE probe, it measures how much detection is
explainable by surface form alone.

\textbf{Within-item AUROC.} Pooled AUROC conflates two abilities: telling
hallucinated \emph{items} from faithful ones (an item-level, topic-correlated
ability) and telling hallucinated \emph{tokens} from faithful ones inside one
answer. Only the second is hallucination localization. The within-item AUROC
isolates it (Fig.~\ref{fig:within}): for every test item containing at least
two tokens of each class ($500/556$ items qualify), we compute AUROC over that
item's tokens and macro-average. Topic, language, and generator are constant
within an item, so a high within-item AUROC cannot arise from topic or surface
confound.

\begin{figure}[t]
\footnotesize
\begin{verbatim}
 pooled AUROC: rank ALL tokens together
   item A (boxing)   item B (anatomy)
   [f f h f h]       [f h h f]
      \____ cross-item topic offset ____/
                (confound)

 within-item AUROC: rank inside each item
   item A: AUROC over {f,f,h,f,h}
   item B: AUROC over {f,h,h,f}
   then macro-average   (topic fixed)

   f = faithful token   h = hallucinated
\end{verbatim}
\caption{Why within-item AUROC removes the topic confound.}
\label{fig:within}
\end{figure}

%==============================================================================
\section{Results}
%==============================================================================

\subsection{Detection and the surface confound}
On the full test split (Table~\ref{tab:detect}), the SAE probe reaches
$0.745 \pm 0.012$. The surface-only baseline reaches $0.698 \pm 0.026$: against
a $0.5$ floor it recovers $0.198/0.245 \approx 80\%$ of the SAE probe's
above-chance discrimination. Surface form is thus a large confound and must be
reported. It does not, however, explain the signal away: per language
(Fig.~\ref{fig:detection}) the SAE probe leads surface by $+0.06$ to $+0.13$ in
all four languages, with non-overlapping error bars. The lead is carried by
genuinely sparse (rare) SAE features: a surface-plus-rare combination
($0.744$) matches the full probe, and rare features alone ($0.739$) nearly do.

\begin{table}[t]
\caption{Detection token-AUROC, full split, 5 seeds.}
\label{tab:detect}
\centering
\begin{tabular}{lc}
\toprule
Feature set & AUROC \\
\midrule
Surface (12 features)      & $0.698 \pm 0.026$ \\
SAE, high-frequency only   & $0.737 \pm 0.016$ \\
SAE, rare only             & $0.739 \pm 0.011$ \\
Surface $+$ SAE rare       & $0.744 \pm 0.010$ \\
SAE, full                  & $\mathbf{0.745 \pm 0.012}$ \\
\bottomrule
\end{tabular}
\end{table}

\begin{figure}[t]
\centering
\includegraphics[width=\columnwidth]{fig_detection.pdf}
\caption{Per-language detection AUROC (full split, 5 seeds): the SAE probe
leads the surface baseline in all four languages, with non-overlapping error
bars. Error bars are $\pm 1$ standard deviation across seeds.}
\label{fig:detection}
\end{figure}

\subsection{Within-item control: the signal is real}
The surface baseline controls form but not topic. The within-item AUROC
controls topic, and it is the decisive result (Fig.~\ref{fig:within-res}). The
SAE probe scores $0.818 \pm 0.011$ within items, well above chance and
$\approx 0.09$ above the surface within-item score of $0.730$, in every
language, with tight non-overlapping bars. Crucially, the within-item score
($0.818$) \emph{exceeds} the pooled score ($0.744$): cross-item topic variation
was depressing the pooled metric, not inflating it. Removing the item-level
nuisance makes the token-level signal stronger. The probe therefore localizes
hallucinated tokens for a reason that is neither surface nor topic.

\begin{figure}[t]
\centering
\includegraphics[width=\columnwidth]{fig_within.pdf}
\caption{Within-item AUROC (topic, language, and generator held constant inside
each answer; $500/556$ items qualify, 5 seeds). The SAE probe beats the surface
within-item control by $\approx 0.09$ in every language and sits above the
pooled cross-item AUROC (dashed), showing the signal is token-level rather than
topic. Error bars are $\pm 1$ standard deviation across seeds.}
\label{fig:within-res}
\end{figure}

\subsection{Span-level IoU}
Strong token discrimination does not convert fully into span overlap on a
base-rate-saturated benchmark, but the gap is smaller than a small-sample read
suggests. On the full split with a per-language threshold
(Table~\ref{tab:iou}), the probe reaches mean char-IoU $0.396 \pm 0.009$
against a FLAG-ALL floor of $0.313 \pm 0.017$ (a lift of $+0.084 \pm 0.023$)
and a FLAG-NONE floor of $\approx 0.05$; the UCSC reference is $\approx 0.55$.
The probe clears FLAG-ALL in all four languages, decisively in es, cs, and en
($+0.10$ to $+0.13$) and to a tie in Mandarin ($0.476$ vs.\ $0.474$), whose
floor is very high because answers there are densely hallucinated and marking
everything is hard to beat. The remaining gap to UCSC we treat as a
span-extraction problem (post-processing, multi-layer features), not a signal
problem.

\begin{table}[t]
\caption{Char-level IoU (full split, 5 seeds, per-language threshold).
Mean $\pm$ std across seeds.}
\label{tab:iou}
\centering
\begin{tabular}{lcc}
\toprule
Language & FLAG-ALL & SAE probe \\
\midrule
es & $0.180 \pm 0.021$ & $\mathbf{0.312 \pm 0.029}$ \\
cs & $0.246 \pm 0.035$ & $\mathbf{0.364 \pm 0.013}$ \\
zh & $0.474 \pm 0.028$ & $\mathbf{0.476 \pm 0.023}$ \\
en & $0.325 \pm 0.040$ & $\mathbf{0.423 \pm 0.024}$ \\
\midrule
Mean & $0.313 \pm 0.017$ & $\mathbf{0.396 \pm 0.009}$ \\
\bottomrule
\end{tabular}
\end{table}

\subsection{Interpretation and ablations}
SAE features match raw residuals (AUROC $0.770$ vs.\ $0.773$): interpretability
costs no accuracy. Inspecting the most heavily weighted features, however,
returns mostly \emph{topic} directions (e.g.\ boxing, anatomy, transport) and
\emph{surface} tokens (sentence-final periods, double newlines, function
words)---which is precisely what motivated the confound controls above. The
token-level signal that survives those controls is distributed across many rare
features rather than concentrated in a few nameable directions.

We tested three elaborations, none of which helped (Table~\ref{tab:neg}).
Routing claims by type (External / Relational / Subjective) before verification
left per-route AUROC flat and routed IoU no better than a flat probe; the
routing itself agrees with a reference only moderately
($\kappa = 0.725$). An internal-vs-external ``gap'' signal added nothing on top
of the probe ($0.769$ vs.\ $0.770$), because it is not measurable on
heterogeneous generations the model did not produce. Proxy retrieval (REFIND)
sat at chance.

\begin{table}[t]
\caption{Negative results / ablations.}
\label{tab:neg}
\centering
\begin{tabular}{lc}
\toprule
Probe & AUROC \\
\midrule
Routing: Relational route & 0.780 \\
Routing: External route    & 0.791 \\
Routing: Subjective route  & 0.708 \\
Confidence ``gap'' alone   & 0.540 \\
SAE $+$ gap                & 0.769 \\
SAE alone                  & 0.770 \\
SAE vs.\ raw residual      & 0.770 / 0.773 \\
REFIND retrieval (CSR)     & 0.502 \\
\bottomrule
\end{tabular}
\end{table}

%==============================================================================
\section{Discussion and Limitations}
%==============================================================================
\textbf{Proxy-model setting.} Gemma-2-2b did not generate the answers it
scores; it reads text produced by other models. That the within-item signal is
strong in this cross-model setting is notable---the model's representation of
arbitrary text encodes token-level factuality---and arguably a strength. It
does, however, mean the internal-vs-external gap method, which requires
self-generation, is out of scope here.

\textbf{What the controls do and do not exclude.} The within-item AUROC
excludes topic, language, and generator confounds and, jointly with the surface
baseline, excludes simple surface form. It does not exclude every conceivable
token-level confound (for instance, entity-token vs.\ function-token effects
not captured by the twelve surface features). The claim is therefore
``robust to surface and topic,'' not ``provably hallucination-only.''

\textbf{IoU.} Competitive span IoU is not yet demonstrated; the contribution is
a validated token-level signal plus the methodology that validates it.

\textbf{Scope.} A single backbone (Gemma-2-2b), single layer (12), and four
languages. Multi-layer features and a larger backbone are natural extensions.

%==============================================================================
\section{Conclusion}
%==============================================================================
A white-box probe over interpretable SAE features detects hallucination at the
token level on multilingual Mu-SHROOM. We show that a trivial surface baseline
recovers most of the naive AUROC---so such numbers must be reported against a
surface control---yet a within-item control, holding topic constant, confirms a
genuine token-level signal ($0.818$ AUROC) in all four languages. The signal is
real; the surface confound is real and under-reported; and the controls that
separate them are the methodological core. Converting the strong token signal
into competitive spans, and testing it under self-generation, are the next
steps.

%==============================================================================
\section*{Reproducibility}
%==============================================================================
All experiments are scripted (feature extraction with cached SAE encodings,
the surface and within-item controls, and the ablations) and reproducible from
the Mu-SHROOM test split with five fixed seeds.

%==============================================================================
\appendices
\section{Per-language AUROC (tabular)}
%==============================================================================
For completeness, Table~\ref{tab:appendix} gives the exact per-language means
and standard deviations underlying Figs.~\ref{fig:detection}
and~\ref{fig:within-res} (full split, 5 seeds).

\begin{table*}[t]
\renewcommand{\arraystretch}{1.3}
\caption{Per-language token-AUROC, mean $\pm$ std across 5 seeds.
``Detection'' is pooled; ``Within-item'' holds topic constant inside each item.}
\label{tab:appendix}
\centering
\begin{tabular}{l cc cc}
\toprule
& \multicolumn{2}{c}{Detection (pooled)} & \multicolumn{2}{c}{Within-item (topic held constant)} \\
\cmidrule(lr){2-3}\cmidrule(lr){4-5}
Language & Surface & SAE & Surface & SAE \\
\midrule
Spanish (es)  & $0.672 \pm 0.014$ & $0.735 \pm 0.017$ & $0.757 \pm 0.018$ & $0.843 \pm 0.008$ \\
Czech (cs)    & $0.631 \pm 0.032$ & $0.759 \pm 0.014$ & $0.733 \pm 0.027$ & $0.831 \pm 0.014$ \\
Mandarin (zh) & $0.603 \pm 0.013$ & $0.696 \pm 0.024$ & $0.696 \pm 0.018$ & $0.781 \pm 0.017$ \\
English (en)  & $0.655 \pm 0.014$ & $0.757 \pm 0.040$ & $0.740 \pm 0.024$ & $0.827 \pm 0.025$ \\
\midrule
Overall       & $0.698 \pm 0.026$ & $0.745 \pm 0.012$ & $0.730 \pm 0.013$ & $0.818 \pm 0.011$ \\
\bottomrule
\end{tabular}
\end{table*}

%==============================================================================
% References -- inline; verify details before submission.
%==============================================================================
\clearpage
\begin{thebibliography}{99}
\bibitem{mushroom} R.~V\'azquez \emph{et al.}, ``SemEval-2025 Task~3:
Mu-SHROOM, the multilingual shared-task on hallucinations,''
\emph{Proc.\ SemEval}, 2025. 
\bibitem{ucsc} J.~Huang \emph{et al.}, ``UCSC at SemEval-2025 Task~3,''
\emph{Proc.\ SemEval}, 2025. 
\bibitem{deloitte} ``Deloitte (Drocks) at SemEval-2025 Task~3,''
\emph{ACL Anthology} 2025.semeval-1.144, 2025. 
\bibitem{orgad} H.~Orgad \emph{et al.}, ``LLMs know more than they show:
on the intrinsic representation of hallucinations,'' \emph{ICLR}, 2025.
\bibitem{ferrando} J.~Ferrando \emph{et al.}, ``Do I know this entity?
Knowledge awareness and hallucinations in language models,''
\emph{ICLR}, 2025.
\bibitem{skeptic} ``On the limits of internal-state probes for hallucination
detection,'' arXiv:2510.09033, 2025.
\bibitem{gemmascope} T.~Lieberum \emph{et al.}, ``Gemma Scope: open sparse
autoencoders everywhere all at once on Gemma~2,'' arXiv:2408.05147, 2024.
\bibitem{gemma2} Gemma Team, ``Gemma~2: improving open language models at a
practical size,'' Google DeepMind, 2024.
\bibitem{refind} D.~Lee and S.~Yu, ``REFIND: retrieval-augmented factuality
hallucination detection,'' arXiv:2502.13622, 2025. 
\bibitem{monosemantic} T.~Bricken \emph{et al.}, ``Towards monosemanticity:
decomposing language models with dictionary learning,'' Anthropic, 2023.
\bibitem{cunningham} H.~Cunningham \emph{et al.}, ``Sparse autoencoders find
highly interpretable features in language models,'' \emph{ICLR}, 2024.
\bibitem{azaria} A.~Azaria and T.~Mitchell, ``The internal state of an LLM
knows when it's lying,'' \emph{Findings of EMNLP}, 2023. 
\end{thebibliography}
\end{document}